% Generated draft: requires mathematical and source review.
\documentclass[11pt]{article}
\usepackage[margin=25mm]{geometry}
\usepackage{fontspec}
\setmainfont{Latin Modern Roman}
\usepackage{amsmath,amssymb,mathtools,unicode-math}
\setmathfont{Latin Modern Math}
\usepackage{enumitem,needspace,xurl,hyperref}
\setlength{\parindent}{0pt}
\setlength{\parskip}{5pt}
\setlength{\emergencystretch}{3em}
\newcommand{\N}{\mathbb{N}}
\newcommand{\Z}{\mathbb{Z}}
\newcommand{\Q}{\mathbb{Q}}
\newcommand{\R}{\mathbb{R}}
\newcommand{\C}{\mathbb{C}}
\newcommand{\F}{\mathbb{F}}
\newcommand{\A}{\mathbb{A}}
\newcommand{\PP}{\mathbb{P}}
\newcommand{\E}{\mathbb{E}}
\newcommand{\eps}{\varepsilon}
\newcommand{\dd}{\,\mathrm{d}}
\begin{document}
\section*{UnsolvedMath ID: 30006042}
\textbf{Source problem code:} catalogue entry
\subsection{Definitions and mathematical statement}
Let \(n\geq4\), and let \(\Omega\) be an \(n\)-dimensional manifold with smooth boundary \(\Sigma=\partial\Omega\) satisfying \(\pi_1(\Omega,\Sigma)=0\), as in the supplied theorem excerpt. The source asks whether its generic regularity and local-solvability conclusions extend to positive Einstein constants. The following describes the intended boundary map and conclusions, subject to the unresolved qualifications stated below.

An Einstein metric with positive Einstein constant is a Riemannian metric \(g\) satisfying
\[
\operatorname{Ric}_g=\lambda g,\qquad \lambda>0,
\]
where \(\lambda\) is constant on \(\Omega\). Write \(g^\top\) for its induced boundary metric, \([g^\top]\) for the boundary conformal class, and \(H_g\) for the boundary mean curvature. Here two boundary metrics \(\gamma_1,\gamma_2\) are conformally equivalent when \(\gamma_2=e^{2u}\gamma_1\) for a real-valued function \(u\) of the required regularity. The mean-curvature normal and normalization must be fixed consistently; the supplied record does not specify them.

To display the question without deciding whether the Einstein constant is fixed, let \(\Lambda\subset(0,\infty)\) denote the permitted constants and put formally
\[
\mathcal E^{k,\alpha}_{\Lambda}(\Omega)
=\{g\text{ of class }C^{k,\alpha}:\operatorname{Ric}_g=\lambda g\text{ for some }\lambda\in\Lambda\}.
\]
The parameters \(k,\alpha\) must belong to the regularity range of the cited theorem, which is not supplied. Let \(\mathcal C^{k,\alpha}(\Sigma)\) denote conformal classes of boundary metrics of that regularity. With the conventional choice of \(\mathcal D^{k+1,\alpha}_{\partial}(\Omega)\) as diffeomorphisms restricting to the identity on \(\Sigma\), the Anderson boundary map has the formal expression
\[
\Pi_{\Lambda}:\mathcal E^{k,\alpha}_{\Lambda}(\Omega)/\mathcal D^{k+1,\alpha}_{\partial}(\Omega)
\longrightarrow
\mathcal C^{k,\alpha}(\Sigma)\times C^{k-1,\alpha}(\Sigma),
\qquad
[g]\longmapsto([g^\top],H_g).
\]
This expression describes the boundary data; identifying it with the precise domain and target in the source requires clarification of the corrupted notation and any additional restrictions.

\textbf{Question, conditional on the intended analytic setting.} Does the positive-constant version of this map have an open dense set of regular points? In the usual differential interpretation, regularity at \([g]\) means that \(D\Pi_{\Lambda}|_{[g]}\) is surjective. Furthermore, for each regular \([g]\), is there a neighborhood \(V\) of \(([g^\top],H_g)\) such that, for every \(([\gamma],\phi)\in V\), there exist \(\varepsilon>0\) and a family of Einstein metrics \((g_s)_{|s|<\varepsilon}\) in the intended domain satisfying
\[
\Pi_{\Lambda}([g_s])=([\gamma],\phi)
\qquad\text{for every }|s|<\varepsilon?
\]
The topology, the precise meaning of regularity, and any required regularity or nontriviality of this family must be those intended by the cited theorem.

The supplied text does not determine whether \(\Lambda=\{\lambda_+\}\) for a prescribed \(\lambda_+>0\), whether the conclusion is required separately for every such \(\lambda_+\), or whether \(\Lambda=(0,\infty)\) with the constant allowed to vary within a family. It also does not specify compactness of \(\Omega\), the full metric-space restrictions, or the precise boundary-data space. Consequently, the displayed formulation records the intended question but does not resolve these missing hypotheses.
\subsection{Short English statement}
For an Einstein metric on a manifold with boundary, Anderson boundary data consist of the conformal class of the induced boundary metric and the boundary mean curvature. The question asks whether the supplied generic regularity and nearby-data solvability conclusions for Ricci-flat and negative Einstein metrics extend to positive Einstein constants, under the stated dimension and relative fundamental-group hypotheses. The exact analytic spaces, treatment of the Einstein constant, and meaning of the asserted family remain unresolved in the supplied excerpt.
\subsection{Sources}
\begin{itemize}
\item Local UnsolvedMath registry, permanent ID \(30006042\), exact source problem code \texttt{catalogue entry}, title ``Einstein Boundary Regularity with Positive Cosmological Constant.'' The supplied complete record includes the question and an extracted Theorem 3. No individual proposer is identified.
\item Supplied source citation: \emph{Mathematical Aspects of General Relativity (2025)}, Oberwolfach Reports, pp.\ 2033--2112, DOI \texttt{10.4171/owr/2024/36}, \url{https://doi.org/10.4171/owr/2024/36}. The supplied literature-source label is ``Oberwolfach Report 36 (2024),'' with the same link. The difference between the citation's 2025 designation and the report's 2024 designation is retained without reconciliation.
\item Literature assessment dated \texttt{2026-08-22}: the supplied curation labels the problem open, says the source explicitly raises the positive-Einstein-constant analogue, and reports that no later resolution was located. This is supplied dated curation, not an independently verified literature assessment. Its listed source is \url{https://doi.org/10.4171/owr/2024/36}.
\item The registry labels the statement \texttt{corrected\_verified} and records a prior check of the official report and a repair of notation. These are supplied verification claims; no independent source visit or mathematical verification is asserted here.
\end{itemize}
\end{document}
